\documentclass[11pt]{article}
\usepackage{amsmath,amssymb}
\usepackage[hidelinks]{hyperref}

\title{Warm-up Matters: Stabilising Small Transformers}
\author{Ada Lovelace}
\date{\today}

\begin{document}
\maketitle

\begin{abstract}
A short learning-rate warm-up removes almost every divergent run when training small transformers, at no cost in final perplexity.
\end{abstract}

\section{Introduction}

Transformers (Vaswani et al. 2017) are known to train unstably when the learning rate starts high.
We ask how much warm-up a small model needs, and whether its shape matters.

\section{Method}

learning rate at step t: eta_t = eta_max * min(1, t/T_w), T_w = warm-up length

loss: L(theta) = -1/N sum_i log p(y_i | x_i) + lambda ||theta||^2

\section{Results}

WikiText-103, 5 seeds each:
none: ppl 31.8, 3/5 diverged
linear 1k: ppl 24.6, 0/5 diverged
linear 4k: ppl 24.1, 0/5 diverged
cosine 4k: ppl 24.3, 0/5 diverged

\section{Conclusion}


\end{document}
